% GENERATED FILE. Edit canonical lessons, then run npm run generate:books.
% canonical-source-hash: fnv1a64:391d8bfee50cb4a4
% canonical-lessons: FA-C366-bahs, FA-C366-bahs-kardan, FA-C366-mozakere, FA-C366-eteraz, FA-C366-mesal-zadan

\chapter{Discussion and debate}
\label{ch:fa-discussion-and-debate}
\hlchaptermodality{366}

\begin{chapteropening}
\textbf{By the end of this chapter:} \emph{I can talk about a discussion or a debate, say that people argued, mention talks between countries, object and give an example.}

Complete the last lesson of chapter 366: I can talk about a discussion or a debate, say that people argued, mention talks between countries, object and give an example.
\end{chapteropening}

\section[\texorpdfstring{\fa{بحث} (bahs)}{بحث (bahs)}]{\fa{بحث} (bahs) — a discussion, a debate}

\label{lesson:FA-C366-bahs}



\begin{quote}
Before the new one: say the Persian for to exaggerate, then the Persian for to justify.
\end{quote}

\subsection*{You'll want to know: \fa{بحث}}
\textbf{\fa{بحث}} (\emph{bahs}) — \textquotedblleft{}a discussion, a debate\textquotedblright{}.

\begin{quote}
\fa{بعد} \fa{از} \fa{جلسه} \fa{یک} \fa{بحث} \fa{طولانی} \fa{داشتیم}.
\end{quote}

\begin{quote}
\emph{ba'd az jalase yek bahs-e tulâni dâshtim.}
\end{quote}

\begin{quote}
After the meeting we had a long discussion.
\end{quote}

Plural \emph{bahshâ}, or formal \emph{mabâhes}. In Persian, \fa{ث} sounds like s.

\begin{cousinweb}
\textbf{\fa{بحث}} was borrowed from Arabic \emph{baḥth} \textquotedblleft{}search, examination, discussion\textquotedblright{} (EtymDB): a debate is a search for the answer.
\end{cousinweb}

\begin{grammarlens}[title={What you've built}]
One more word for telling, explaining and describing at B1.
\end{grammarlens}

\subsection*{Your turn}
\begin{itemize}
\raggedright
  \item \emph{Say it:} \emph{bahs}
  \item \emph{Say it:} \emph{bahs}, once more
  \item \emph{Say it:} a sentence of your own with \emph{bahs}
  \item \emph{Recall:} say the Persian for to exaggerate, then the Persian for to justify, then say \emph{bahs} again
\end{itemize}

\begin{tcolorbox}[breakable,colback=teal!4,colframe=teal!35!black,title={Before you move on}]
What does \fa{بحث} mean? (\textquotedblleft{}A discussion, a debate\textquotedblright{}.) Say it once more.
\end{tcolorbox}
\section[\texorpdfstring{\fa{بحث} \fa{کردن} (bahs kardan)}{بحث کردن (bahs kardan)}]{\fa{بحث} \fa{کردن} (bahs kardan) — to discuss, to argue}

\label{lesson:FA-C366-bahs-kardan}



\begin{quote}
Before the new one: say the Persian for to justify, then the Persian for a discussion.
\end{quote}

\subsection*{You'll want to know: \fa{بحث} \fa{کردن}}
\textbf{\fa{بحث} \fa{کردن}} (\emph{bahs kardan}) — \textquotedblleft{}to discuss, to argue\textquotedblright{}.

\begin{quote}
\fa{آنها} \fa{درباره} \fa{پول} \fa{بحث} \fa{میکنند}.
\end{quote}

\begin{quote}
\emph{ânhâ darbâre-ye pul bahs mikonand.}
\end{quote}

\begin{quote}
They are arguing about money.
\end{quote}

\emph{Bahs kardan darbâre-ye} … — to argue about … ; with someone: \emph{bâ kasi}.

\begin{cousinweb}
\textbf{\fa{بحث} \fa{کردن}} joins \emph{bahs}, from Arabic \emph{baḥth} \textquotedblleft{}discussion\textquotedblright{} (EtymDB), to the inherited Persian verb \emph{kardan} \textquotedblleft{}to do\textquotedblright{}.
\end{cousinweb}

\begin{grammarlens}[title={What you've built}]
One more word for telling, explaining and describing at B1.
\end{grammarlens}

\subsection*{Your turn}
\begin{itemize}
\raggedright
  \item \emph{Say it:} \emph{bahs kardan}
  \item \emph{Say it:} \emph{bahs kardan}, once more
  \item \emph{Say it:} a sentence of your own with \emph{bahs kardan}
  \item \emph{Recall:} say the Persian for to justify, then the Persian for a discussion, then say \emph{bahs kardan} again
\end{itemize}

\begin{tcolorbox}[breakable,colback=teal!4,colframe=teal!35!black,title={Before you move on}]
What does \fa{بحث} \fa{کردن} mean? (\textquotedblleft{}To discuss, to argue\textquotedblright{}.) Say it once more.
\end{tcolorbox}
\section[\texorpdfstring{\fa{مذاکره} (mozâkere)}{مذاکره (mozâkere)}]{\fa{مذاکره} (mozâkere) — a negotiation, talks}

\label{lesson:FA-C366-mozakere}



\begin{quote}
Before the new one: say the Persian for a discussion, then the Persian for to discuss.
\end{quote}

\subsection*{You'll want to know: \fa{مذاکره}}
\textbf{\fa{مذاکره}} (\emph{mozâkere}) — \textquotedblleft{}a negotiation, talks\textquotedblright{}.

\begin{quote}
\fa{مذاکره} \fa{دو} \fa{کشور} \fa{سه} \fa{روز} \fa{طول} \fa{کشید}.
\end{quote}

\begin{quote}
\emph{mozâkere-ye do keshvar se ruz tul keshid.}
\end{quote}

\begin{quote}
The talks between the two countries lasted three days.
\end{quote}

Formal talks are often plural: \emph{mozâkerât}. \emph{Mozâkere kardan} — to negotiate.

\begin{cousinweb}
\textbf{\fa{مذاکره}} comes from Arabic \emph{mudhākara} (EtymDB), a mutual verbal noun usually explained as reminding one another; in Persian \fa{ذ} sounds like z.
\end{cousinweb}

\begin{grammarlens}[title={What you've built}]
One more word for telling, explaining and describing at B1.
\end{grammarlens}

\subsection*{Your turn}
\begin{itemize}
\raggedright
  \item \emph{Say it:} \emph{mozâkere}
  \item \emph{Say it:} \emph{mozâkere}, once more
  \item \emph{Say it:} a sentence of your own with \emph{mozâkere}
  \item \emph{Recall:} say the Persian for a discussion, then the Persian for to discuss, then say \emph{mozâkere} again
\end{itemize}

\begin{tcolorbox}[breakable,colback=teal!4,colframe=teal!35!black,title={Before you move on}]
What does \fa{مذاکره} mean? (\textquotedblleft{}A negotiation, talks\textquotedblright{}.) Say it once more.
\end{tcolorbox}
\section[\texorpdfstring{\fa{اعتراض} (e'terâz)}{اعتراض (e'terâz)}]{\fa{اعتراض} (e'terâz) — an objection, a protest}

\label{lesson:FA-C366-eteraz}



\begin{quote}
Before the new one: say the Persian for to discuss, then the Persian for a negotiation.
\end{quote}

\subsection*{You'll want to know: \fa{اعتراض}}
\textbf{\fa{اعتراض}} (\emph{e'terâz}) — \textquotedblleft{}an objection, a protest\textquotedblright{}.

\begin{quote}
\fa{اعتراض} \fa{او} \fa{فایده} \fa{نداشت}.
\end{quote}

\begin{quote}
\emph{e'terâz-e u fâyede nadâsht.}
\end{quote}

\begin{quote}
His objection was no use.
\end{quote}

Plural \emph{e'terâzât}. \emph{Be chizi e'terâz kardan} — to object to something.

\begin{cousinweb}
\textbf{\fa{اعتراض}} was borrowed from Arabic \emph{i'tirāḍ} \textquotedblleft{}objection\textquotedblright{} (EtymDB), usually explained as standing across someone's way.
\end{cousinweb}

\begin{grammarlens}[title={What you've built}]
One more word for telling, explaining and describing at B1.
\end{grammarlens}

\subsection*{Your turn}
\begin{itemize}
\raggedright
  \item \emph{Say it:} \emph{e'terâz}
  \item \emph{Say it:} \emph{e'terâz}, once more
  \item \emph{Say it:} a sentence of your own with \emph{e'terâz}
  \item \emph{Recall:} say the Persian for to discuss, then the Persian for a negotiation, then say \emph{e'terâz} again
\end{itemize}

\begin{tcolorbox}[breakable,colback=teal!4,colframe=teal!35!black,title={Before you move on}]
What does \fa{اعتراض} mean? (\textquotedblleft{}An objection, a protest\textquotedblright{}.) Say it once more.
\end{tcolorbox}
\section[\texorpdfstring{\fa{مثال} \fa{زدن} (mesâl zadan)}{مثال زدن (mesâl zadan)}]{\fa{مثال} \fa{زدن} (mesâl zadan) — to give an example}

\label{lesson:FA-C366-mesal-zadan}



\begin{quote}
Before the new one: say the Persian for a negotiation, then the Persian for an objection.
\end{quote}

\subsection*{You'll want to know: \fa{مثال} \fa{زدن}}
\textbf{\fa{مثال} \fa{زدن}} (\emph{mesâl zadan}) — \textquotedblleft{}to give an example\textquotedblright{}.

\begin{quote}
\fa{معلم} \fa{برای} \fa{هر} \fa{کلمه} \fa{یک} \fa{مثال} \fa{زد}.
\end{quote}

\begin{quote}
\emph{mo'allem barâye har kalame yek mesâl zad.}
\end{quote}

\begin{quote}
The teacher gave an example for each word.
\end{quote}

\emph{Mesâl} \textquotedblleft{}example\textquotedblright{} + \emph{zadan} \textquotedblleft{}to strike\textquotedblright{} (present stem \emph{zan-}). \emph{Masalan} \textquotedblleft{}for example\textquotedblright{} is the adverb.

\begin{cousinweb}
\textbf{\fa{مثال}} is usually traced to Arabic \emph{mithāl} \textquotedblleft{}example, likeness\textquotedblright{}; \emph{zadan} \textquotedblleft{}to strike\textquotedblright{} is the inherited Persian verb that also builds \emph{harf zadan} \textquotedblleft{}to talk\textquotedblright{}.
\end{cousinweb}

\begin{grammarlens}[title={What you've built}]
One more word for telling, explaining and describing at B1.
\end{grammarlens}

\subsection*{Your turn}
\begin{itemize}
\raggedright
  \item \emph{Say it:} \emph{mesâl zadan}
  \item \emph{Say it:} \emph{mesâl zadan}, once more
  \item \emph{Say it:} a sentence of your own with \emph{mesâl zadan}
  \item \emph{Recall:} say the Persian for a negotiation, then the Persian for an objection, then say \emph{mesâl zadan} again
\end{itemize}

\begin{tcolorbox}[breakable,colback=teal!4,colframe=teal!35!black,title={Before you move on}]
What does \fa{مثال} \fa{زدن} mean? (\textquotedblleft{}To give an example\textquotedblright{}.) Say it once more.
\end{tcolorbox}
